\section{Setting: the price of first place}
\label{sec:setting}

\paragraph{Leaderboard and aggregation.}
A leaderboard is a score matrix $S\in\mathbb{R}^{N\times m}$ of $N$ models on $m$ tasks, each task $j$ with a ceiling $c_j$ (e.g.\ 100 points). An aggregation rule $A$ maps the matrix to a score or a ranking of the models. We study six rule groups: the mean (raw, and its $z$-scored, geometric and $p$-mean variants as one family), the median and other quantiles, the minimum, the mean rank (Borda), which with ties counted as half is the mean win rate, the plurality of task wins, and Copeland's pairwise-majority score.

\paragraph{Targeted investment.}
A developer picks a target model $t$ and a boost $\delta\in\mathbb{R}^m_{\ge 0}$; the target's score on task $j$ becomes $\min(S_{tj}+\delta_j,\,c_j)$ and every other score is unchanged. The boost stands for training on data related to the task; we train nothing, a boost is the \emph{input} whose cost we measure, and \Cref{sec:calibration} checks its scale against observed fine-tunes. The developer is static and acts once, rivals do not respond, and boosts on different tasks do not interact (\Cref{sec:discussion}). Per-task boosts are capped by a \emph{boost budget} $B$: $\delta_j\le B$. We report price as a function of $B$ and the smallest feasible budget $B^\ast_A(t)$ at which first place under rule $A$ is attainable. ``First'' means ahead of every rival by a margin $\varepsilon>0$ (ties lose).

\paragraph{Price: breadth, depth, feasibility.}
The investment has two coordinates. Its \emph{breadth} is the number of tasks touched, $k=\lVert\delta\rVert_0$, and its \emph{depth} is the total boost $\lVert\delta\rVert_1$, measured in points of the board's own scale (robustness: in standard deviations of the task, as a fraction of headroom, and under a convex cost; \Cref{sec:robustness}). The \emph{breadth--depth frontier} $C_A(k)$ is the least depth of a plan that makes $t$ first under $A$ using at most $k$ tasks; it is infinite where no plan is feasible under $B$ and the ceilings. The \emph{targeting multiplier} is the depth of the cheapest untargeted plan (the same boost on every task) divided by the depth of the cheapest targeted plan. Across boards with different $m$ we use the breadth fraction $k/m$, the per-task depth $\lVert\delta\rVert_1/m$ and $B$ in task standard deviations.

\paragraph{Incumbents.}
Rules crown different leaders, so a lower price under one rule can mean only that its leader is weaker. We therefore report the price against each rule's own leader and against a common incumbent set, the top models under the board's official rule.
